\bookchapter{บทที่ 18: การสำรองข้อมูลและการย้ายข้อมูล}

\chapterintro{ขอเล่าเรื่องเพื่อนบ้านคนหนึ่ง มือถือตกลงบ่อน้ำพองออกมาก็เปิดไม่ติดแล้ว สิ่งที่เธอเสียใจที่สุดไม่ใช่ตัวเครื่อง แต่คือรูปสามปีที่ไม่เคยสำรองไว้เลย — บันทึกทั้งหมดตั้งแต่ลูกเกิดจนเดินได้ ฟังจบใจฉันไม่สบายมาก เพราะเรื่องนี้เลี่ยงได้ทั้งหมด บทนี้เราจะจริงจังใส่กุญแจด้านความปลอดภัยให้ข้อมูลของคุณ}

\booksection{หนึ่ง เข้าใจก่อน: ในมือถือคุณ อะไรหายแล้วหายจริง}

“สถานะการมีชีวิต” ของข้อมูลแต่ละประเภทต่างกัน ต้องแยกให้ออก:

\begin{itemize}
  \item \textbf{เก็บในเครื่อง หายแล้วหายจริง}: รูปวิดีโอ (ถ้าไม่ได้เปิดซิงก์คลาวด์) โน้ต SMS ประวัติโทรออก ไฟล์ที่ดาวน์โหลด
  \item \textbf{เก็บบนเซิร์ฟเวอร์แอป เปลี่ยนเครื่องล็อกอินใหม่ก็ได้คืน}: รายชื่อเพื่อน WeChat คำสั่งซื้อ Taobao ประวัติรายการในแอปธนาคาร ของพวกนี้เก็บบนเซิร์ฟเวอร์ของบริษัทที่เกี่ยวข้อง ไม่ได้เก็บในมือถือคุณ
  \item \textbf{อยู่เขตสีเทา ต้องจัดการพิเศษ}: ประวัติแชท WeChat (แม้เซิร์ฟเวอร์ WeChat เก็บไว้ แต่คุณดาวน์โหลดจากเซิร์ฟเวอร์เองไม่ได้) อุปกรณ์ยืนยันที่ผูกกับแอปธนาคาร กุญแจในแอปยืนยันสองชั้น (แบบ Google Authenticator)
\end{itemize}

ดังนั้นสิ่งที่คุณต้องคุ้มครองเป็นพิเศษ คือประเภทหนึ่งกับสาม

\booksection{สอง สำรองทั้งเครื่อง: จบทุกอย่างในรอบเดียว}

สำรองทั้งเครื่องเหมือนถ่าย “สแนปช็อต” เต็มใบให้มือถือ — การตั้งค่าระบบปัจจุบัน ผังหน้าจอหลัก ข้อมูลแอป รูปภาพ สมุดโทรศัพท์ แพ็กหมด

\booksubsection{สูตรหนึ่ง: สำรองในเครื่องคอม (น่าเชื่อถือที่สุด แนะนำคนที่มีคอม)}

\begin{itemize}
  \item \textbf{iPhone}: ต่อมือถือเข้าคอมด้วยสายข้อมูล เปิด iTunes (Windows รุ่นเก่า) หรือ Finder (Mac) ที่หน้าอุปกรณ์เลือก "สำรองข้อมูลทันที" อย่าลืมติ๊ก "เข้ารหัสการสำรองข้อมูล" ไม่งั้นรหัสผ่านกับข้อมูลสุขภาพจะไม่ถูกสำรองไปด้วย
  \item \textbf{แอนดรอยด์}: แต่ละแบรนด์มีซอฟต์แวร์ฝั่งคอมของตัวเอง เช่น Huawei HiSuite, Xiaomi Assistant, Samsung Smart Switch ต่อคอมแล้วทำตามคำแนะนำก็จบ
  \item ข้อดีของการสำรองบนคอมคือไม่ต้องซื้อพื้นที่คลาวด์ และคืนข้อมูลเร็วกว่าคลาวด์เยอะ ข้อเสียคือต้องจำได้ว่าต่อสายสำรองเป็นระยะ — คนส่วนใหญ่ทำไม่ได้จุดนี้
\end{itemize}

\booksubsection{สูตรสอง: สำรองคลาวด์ (สบายใจที่สุด แนะนำทุกคน)}

\begin{itemize}
  \item \textbf{iPhone}: ตั้งค่า ชื่อของคุณ iCloud iCloud Backup เปิด เครื่องจะสำรองเองตอนต่อไฟ ต่อ Wi-Fi ล็อกจออยู่ ฟรี 5GB แทบไม่พอ แนะนำจ่าย 6 หยวนอัปเกรดเป็น 50GB
  \item \textbf{แอนดรอยด์}: คลาวด์ของทุกแบรนด์มีฟังก์ชันสำรอง Huawei ไปค้นหา "การสำรองข้อมูลคลาวด์" ในตั้งค่า Xiaomi ไป "บริการคลาวด์ Xiaomi" OPPO/vivo ก็คล้ายกัน นอกจากนี้ยังมีการสำรองในตัวของ Google (ตั้งค่า ระบบ การสำรองข้อมูล) เปิดมันด้วยเลย
  \item ข้อดีใหญ่สุดของสำรองคลาวด์คือ "คุณไม่ต้องเป็นห่วง" จัดการเองหมด ข้อเสียใหญ่สุดคือตอนคืนต้องดาวน์โหลดนาน และพื้นที่ฟรีมักไม่พอ
\end{itemize}

\booksection{สาม กลยุทธ์สำรองแยกประเภท: ของสำคัญเก็บแยกอีกชุด}

สำรองทั้งเครื่องคือตาข่ายด้านล่าง แต่บางอย่างสำคัญเกินกว่าจะไม่เก็บแยกอีกชุด:

\booksubsection{รูปภาพและวิดีโอ (ทรัพย์สินดิจิทัลล้ำค่าที่สุด)}

\begin{itemize}
  \item \textbf{สูตรประกันสองชั้น}: เปิดซิงก์อัตโนมัติ iCloud Photos/Google Photos พร้อมกันกับทุกสองสามเดือนคัดลอกรูปไปเก็บบนคอมหรือฮาร์ดดิสก์พกพาด้วยมือ คลาวด์หนึ่งชุด เครื่องหนึ่งชุด สำรองให้กันเอง
  \item \textbf{สูตรไม่อยากเสียเงิน}: Google Photos ให้ฟรี 15GB ตั้ง "คุณภาพบีบอัด" ประหยัดพื้นที่ได้ หรือทุกเดือนใช้สายข้อมูลคัดลอกรูปทั้งหมดไปคอม
\end{itemize}

\booksubsection{สมุดโทรศัพท์ (มองข้ามง่ายแต่สำคัญสุดๆ)}

\begin{itemize}
  \item สมุดโทรศัพท์ iPhone ซิงก์กับ iCloud เป็นค่าเริ่มต้น เปิดซิงก์ไว้ก็ไม่ต้องสำรองเพิ่ม ถ้าอยากปลอดภัยอีกขั้น ไป iCloud.com ส่งออกไฟล์ vCard เก็บบนคอม
  \item สมุดโทรศัพท์แอนดรอยด์มักซิงก์กับ Google หรือคลาวด์แบรนด์ ไปที่การตั้งค่าแอปรายชื่อโทรศัพท์ส่งออกไฟล์ .vcf ก็เป็นสำรองฉบับสมบูรณ์หนึ่งชุด
\end{itemize}

\booksubsection{ประวัติแชท WeChat (จัดการแยก สำคัญสุดของสำคัญ)}

ตัวนี้ต้องแยกมาพูดเป็นพิเศษ เพราะการสำรองทั้งเครื่องมักไม่ครบเรื่องประวัติแชท WeChat และเครื่องมือย้ายเครื่องเอามันไปไม่ได้เลย

\begin{itemize}
  \item \textbf{วิธีหนึ่ง: ย้ายไปมือถือใหม่}. WeChat ฉัน ตั้งค่า ทั่วไป การย้ายและสำรองประวัติแชท ย้าย ย้ายไปยังโทรศัพท์/แท็บเล็ต เครื่องเก่าใหม่ต่อ Wi-Fi เดียวกัน สแกนคิวอาร์ก็ส่งได้ ระวัง: ก่อนย้ายเสร็จ อย่าล็อกเอาต์ WeChat จากเครื่องเก่า
  \item \textbf{วิธีสอง: สำรองไปคอม}. ติดตั้ง WeChat PC บนคอม สามขีดมุมซ้ายล่าง การย้ายและสำรอง การสำรองและกู้คืน สำรองประวัติแชทไปคอม หลังเปลี่ยนเครื่องคืนลงมือถือใหม่ได้
  \item \textbf{วิธีสาม: สำรองไปคลาวด์}. WeChat รองรับสำรองไป iCloud หรือพื้นที่คลาวด์มือถือด้วย แต่ระวัง: การสำรองแบบนี้คืนได้เฉพาะบนมือถือเครื่องปัจจุบัน ข้ามอุปกรณ์ไม่ได้ เปลี่ยนเครื่องต้องใช้สองวิธีแรก
\end{itemize}

\booksection{สี่ คู่มือเปลี่ยนเครื่อง: จากเครื่องเก่าสู่เครื่องใหม่}

ได้เครื่องใหม่เป็นเรื่องดี แต่ถ้าข้อมูลย้ายไม่ไปก็กลายเป็นปวดหัว จัดการตามสามสถานการณ์:

\booksubsection{อัปเกรดแบรนด์เดิม (ง่ายที่สุด)}

\begin{itemize}
  \item \textbf{iPhone iPhone เครื่องใหม่}: ตอนเปิดเครื่องเปิดใช้งาน วางเครื่องเก่าไว้เคียงข้าง จะมี "เริ่มต้นด่วน" โผล่ขึ้นมาเอง ใช้เครื่องเก่าสแกนกราฟิกบนเครื่องใหม่ก็จบ ข้อมูล การตั้งค่า ผังแอป ย้ายไปครบแบบไม่เปลี่ยน รวมถึงรหัส Wi-Fi ก็เอาไปให้อัตโนมัติ ใช้เวลาราว 20-30 นาที
  \item \textbf{แอนดรอยด์แบรนด์เดียวกัน}: ใช้เครื่องมือย้ายเครื่องของแบรนด์ Huawei ใช้ "Phone Clone" Xiaomi ใช้ "Mi Mover" OPPO ใช้ "โอนย้ายข้อมูล" vivo ใช้ "EasyShare" ขั้นตอนใกล้เคียงกันหมด: ติดตั้งแอปเดียวกันบนเครื่องเก่าใหม่ สแกนเชื่อมต่อ ติ๊กของที่จะย้าย เริ่มส่ง
\end{itemize}

\booksubsection{เปลี่ยนข้ามแบรนด์ (แอนดรอยด์สู่แอนดรอยด์ คนละแบรนด์)}

ใช้เครื่องมือย้ายเครื่องภายนอกหรือแอปย้ายเครื่องของแบรนด์เก่า (แอปย้ายเครื่องของแบรนด์ส่วนใหญ่รองรับย้ายไปแอนดรอยด์แบรนด์อื่นด้วย) ย้ายได้: สมุดโทรศัพท์ รูป SMS ประวัติโทรออก แอปบางตัว ย้ายไม่ได้: สถานะล็อกอินในแอป (ต้องล็อกอินใหม่) ประวัติแชท WeChat (ใช้ฟังก์ชันย้ายของ WeChat เอง)

\booksubsection{เปลี่ยนข้ามแพลตฟอร์ม (แอนดรอยด์สู่ iPhone หรือกลับกัน)}

\begin{itemize}
  \item \textbf{แอนดรอยด์ iPhone}: แอปเปิลมีแอป "Move to iOS" ดาวน์โหลดลงมือถือแอนดรอยด์ แล้วตอนเปิดใช้งาน iPhone ใหม่เลือก "ย้ายข้อมูลจากอุปกรณ์ Android" ย้ายสมุดโทรศัพท์ รูป SMS ข้อมูลพื้นฐานได้ แต่ข้อมูลในแอปย้ายไม่ได้ ประวัติแชท WeChat ยังต้องใช้ฟังก์ชันย้ายของ WeChat เอง
  \item \textbf{iPhone แอนดรอยด์}: แต่ละแบรนด์แอนดรอยด์มีเครื่องมือย้ายที่สอดคล้อง แต่พูดตรงๆ ประสบการณ์ย้ายข้ามแพลตฟอร์มสู้แพลตฟอร์มเดียวกันไม่ได้ สำคัญที่สุดทำแยกไว้ก่อน: ส่งออกสมุดโทรศัพท์ สำรองรูป ย้ายประวัติแชท WeChat ที่เหลือค่อยๆ ทำ
\end{itemize}

\booksection{ห้า กฎทองของการสำรอง: หลัก 3-2-1}

นี่คือหลักการคลาสสิกในวงการความปลอดภัยข้อมูล จำง่าย:

\begin{itemize}
  \item \textbf{3}: ข้อมูลสำคัญเก็บไว้อย่างน้อยสามชุด (ต้นฉบับหนึ่ง สำรองสองชุด)
  \item \textbf{2}: สำรองลงสื่อจัดเก็บสองแบบที่ต่างกัน (เช่นคลาวด์หนึ่ง ฮาร์ดดิสก์พกพาหนึ่ง)
  \item \textbf{1}: อย่างน้อยหนึ่งชุดสำรองอยู่คนละสถานที่กายภาพ (กันไฟไหม้ โจรกรรม พร้อมกันทำลายข้อมูลทั้งหมด)
\end{itemize}

สำหรับคนทั่วไป ทำได้ถึง “คลาวด์หนึ่ง คอมหนึ่ง” ก็ดีมากแล้ว ไม่ต้องซับซ้อนเกิน หัวใจคือคุณทำมันจริงๆ

\begin{tcbwarn}{ประวัติแชท WeChat ระวังเป็นพิเศษ}
ประวัติแชท WeChat ย้ายผ่านเครื่องมือย้ายเครื่องใดๆ ไม่ได้! มันอยู่นอกขอบเขตครอบคลุมของการสำรองทั้งเครื่อง ต้องใช้ฟังก์ชัน “การย้ายประวัติแชท” ในตัวของ WeChat หรือสำรองไป WeChat PC หลายคนเปลี่ยนเครื่องใหม่เสร็จแล้วค่อยเจอว่าประวัติแชทหายหมด ตอนนั้นสายไปแล้ว — เครื่องเก่าอาจถูกรีเซ็ตไปเรียบร้อย

\end{tcbwarn}

\begin{tcbpractice}{ฝึกหน่อย: ทำการซ้อมสำรองข้อมูลฉบับเต็มหนึ่งรอบ}
\begin{enumerate}
  \item ตรวจว่าเปิดการสำรองคลาวด์ iCloud หรือคลาวด์แบรนด์อยู่ไหม ยืนยันเวลาสำรองล่าสุด
  \item สำรองรูปด้วยมือ: ต่อคอมด้วยสายข้อมูล คัดลอกโฟลเดอร์ DCIM และ Pictures ไปคอม
  \item สำรองสมุดโทรศัพท์: ผู้ใช้ iPhone ไป iCloud.com ส่งออก vCard ผู้ใช้แอนดรอยด์ส่งออกไฟล์ .vcf ในการตั้งค่าแอปรายชื่อ
  \item สำรองประวัติแชท WeChat: เปิด WeChat ฉัน ตั้งค่า ทั่วไป การย้ายและสำรองประวัติแชท สำรองไปคอม
  \item คัดลอกไฟล์ vcf สมุดโทรศัพท์กับโฟลเดอร์รูปที่ส่งออกไปเก็บอีกชุดบนฮาร์ดดิสก์พกพาหรือแฟลชไดรฟ์ จบเรื่อง "หนึ่งชุดออฟไลน์อยู่ที่อื่น" ในหลัก 3-2-1
\end{enumerate}

\end{tcbpractice}
